\chapter{Testowanie (Karina Wołoszyn)}
\label{chap:tests}

Testowanie jest integralną częścią procesu tworzenia oprogramowania. Dzięki niemu, jesteśmy w stanie wyłapać defekty w kodzie i błędy w aplikacji, na które użytkownicy końcowi mogliby natrafić podczas korzystania z niej. 

Testowanie pozwala zwiększyć satysfakcję i zaufanie użytkowników do aplikacji i minimalizuje ryzyko ponoszenia większych nakładów czasowych i pracy związanych z naprawą błędów. 

\section{Podejście do testowania}

Podczas pracy nad projektem, zdecydowaliśmy się skupić na przeprowadzaniu testów manualnych, które pozwalają na dokładne sprawdzenie funkcjonalności od perspektywy użytkownika. Dzięki temu mogliśmy mieć pewność, że żadne kroki nie zostaną pominięte, a aplikacja działa zgodnie z oczekiwaniami. Jednak testy manualne mogę być czasochłone i podatne na błędy ludzkie, dlatego też niektóre części naszego kodu zostały również objęte testami automatycznymi.


\section{Zapewnienie jakości}

Jednym z efektów testowania jest zapewnienie jakości oprogramowania - QA (ang. \textit{Quality Assurance}). QA to proces, który ma na celu zapewnienie, że oprogramowanie spełnia określone standardy jakości i działa zgodnie z oczekiwaniami użytkowników. QA zawiera 5 rodzaji aktywności \cite{baresipezze2006}:
- planowanie i monitorowanie
- weryfikacja specyfikacji wymagań
- tworzenie przypadków użycia
- walidacja
- ulepszenie procesu

\section{Testy manualne}

Testy manualne przebiegały na każdym etapie rozwoju aplikacji.


\section{Testy jednostkowe}

Opierając się na piramidzie testowania, stworzyliśmy wiele testów jednostkowych. Testy te, pomimo małego pokrycia kodu, pozwoliły na upewnienie się, czy kluczowe funkcje spełniają swoje zadania poprawnie, a w przypadku wykrycia niezgodności łatwej ich lokalizacji.

\section{Testy e2e}

Testy e2e chodź nie obejmują całego zakresu aplikacji, pozwalają na sprawdzenie najważniejszych ścieżek użytkownika, takich jak rejestracja, logowanie czy rozegranie rozgrywki.



